% ============================================================
\chapter{Bayesian Estimation and Loss Functions}
\label{ch:estimation}
% ============================================================

In frequentist statistics, an estimator is ``good'' if it is unbiased and has minimum variance. In Bayesian statistics, the question is framed differently: which summary of the posterior distribution should one choose, and at what cost? The answer depends on what one considers a serious error. Is being far off proportionally worse than being slightly off (quadratic loss), or is any error beyond a threshold equally unacceptable (0-1 loss)? Abraham Wald, in his foundational work on decision theory in the 1940s, showed that this question of the \emph{loss function} is unavoidable: every estimator is, consciously or not, the solution to a decision problem.

This chapter explores the link between estimation and decision, and shows how each loss function leads to a different optimal estimator: the posterior mean, the posterior median, the posterior mode.

\begin{intuition}[title=\textbf{Central idea}]
Bayesian estimation is not limited to a single point estimator: it
delivers the entire posterior distribution. Nevertheless, when a point
summary is needed, the optimal estimator depends on the chosen
\textbf{loss function}.
\end{intuition}

% ------------------------------------------------------------
\section{Bayesian decision theory}
% ------------------------------------------------------------

\begin{definition}[Decision problem]
A \textbf{decision problem} is a triple $(\Theta, \mathcal{A}, L)$
where:
\begin{itemize}
  \item $\Theta$ is the parameter space,
  \item $\mathcal{A}$ is the action (decision) space,
  \item $L : \Theta \times \mathcal{A} \to \R^+$ is the loss function.
\end{itemize}
\end{definition}

\begin{definition}[Posterior risk]
The \textbf{posterior risk} (or posterior expected loss) of an action
$a \in \mathcal{A}$ is:
\[
  \rho(\pi, a \mid x) = \E_{\theta \mid x}[L(\theta, a)]
  = \int_\Theta L(\theta, a)\,\pi(\theta \mid x)\,d\theta.
\]
\end{definition}

\begin{definition}[Bayes estimator]
The \textbf{Bayes estimator} under loss $L$ is the action minimizing
the posterior risk:
\[
  \hat{\theta}_B(x) = \argmin_{a \in \mathcal{A}} \rho(\pi, a \mid x).
\]
\end{definition}

\begin{definition}[Integrated Bayes risk]
The \textbf{Bayes risk} of an estimator $\delta(x)$ is:
\[
  r(\pi, \delta) = \E_x\left[\rho(\pi, \delta(x) \mid x)\right]
  = \int \int L(\theta, \delta(x))\,f(x \mid \theta)\,\pi(\theta)\,dx\,d\theta.
\]
\end{definition}

\begin{theorem}[Optimality of the Bayes estimator]
The Bayes estimator $\hat{\theta}_B$ minimizes the Bayes risk
$r(\pi, \delta)$ among all estimators.
\end{theorem}

\begin{proof}
The Bayes risk can be written as:
\[
  r(\pi, \delta) = \int \rho(\pi, \delta(x) \mid x)\,m(x)\,dx,
\]
where $m(x)$ is the marginal likelihood. Since $m(x) \geq 0$,
minimizing $r(\pi, \delta)$ amounts to minimizing
$\rho(\pi, \delta(x) \mid x)$ for each $x$, yielding $\hat{\theta}_B(x)$.
\end{proof}

% ------------------------------------------------------------
\section{Classical loss functions}
% ------------------------------------------------------------

\begin{keyformulas}[title=\textbf{Bayes estimators by loss function}]
\begin{center}
\renewcommand{\arraystretch}{1.5}
\begin{tabular}{p{4.5cm} p{4cm} p{5cm}}
\toprule
\textbf{Loss function} & \textbf{Expression} & \textbf{Bayes estimator} \\
\midrule
Quadratic & $L(\theta, a) = (\theta - a)^2$ &
  $\E[\theta \mid x]$ (posterior mean) \\
Absolute value & $L(\theta, a) = \abs{\theta - a}$ &
  $\text{Med}[\theta \mid x]$ (posterior median) \\
0--1 (all-or-nothing) & $L(\theta, a) = \mathbf{1}(\abs{\theta - a} > \varepsilon)$ &
  Mode$[\theta \mid x]$ (MAP) \\
LINEX & $L(\theta, a) = e^{c(\theta-a)} - c(\theta-a) - 1$ &
  $\frac{1}{c}\log \E[e^{c\theta} \mid x]$ \\
Weighted quadratic & $L(\theta, a) = w(\theta)(\theta - a)^2$ &
  $\frac{\E[w(\theta)\theta \mid x]}{\E[w(\theta) \mid x]}$ \\
\bottomrule
\end{tabular}
\end{center}
\end{keyformulas}

\subsection{Quadratic loss}

\begin{theorem}[Bayes estimator under quadratic loss]
\label{thm:bayes_quad}
Under quadratic loss $L(\theta, a) = (\theta - a)^2$, the Bayes
estimator is the posterior mean:
\[
  \hat{\theta}_B(x) = \E[\theta \mid x].
\]
\end{theorem}

\begin{proof}
\begin{align*}
  \frac{\partial}{\partial a} \E[(\theta - a)^2 \mid x]
  &= -2\,\E[\theta - a \mid x]
  = -2\big(\E[\theta \mid x] - a\big).
\end{align*}
This vanishes at $a = \E[\theta \mid x]$. The second derivative is
$2 > 0$, confirming a minimum.
\end{proof}

\begin{corollary}
The posterior risk under quadratic loss is the posterior variance:
\[
  \rho(\pi, \hat{\theta}_B \mid x) = \Var[\theta \mid x].
\]
\end{corollary}

\subsection{Absolute value loss}

\begin{theorem}[Bayes estimator under absolute loss]
Under $L(\theta, a) = \abs{\theta - a}$, the Bayes estimator is the
posterior median.
\end{theorem}

\begin{proof}
The generalized derivative of the posterior risk is:
\[
  \frac{\partial}{\partial a} \E[\abs{\theta - a} \mid x]
  = \PP(\theta < a \mid x) - \PP(\theta > a \mid x),
\]
which vanishes when $\PP(\theta < a \mid x) = 1/2$, i.e., $a$ is the
median of $\pi(\theta \mid x)$.
\end{proof}

\subsection{Maximum a posteriori (MAP)}

\begin{definition}[MAP estimator]
The \textbf{maximum a posteriori} (MAP) estimator is:
\[
  \hat{\theta}_{\text{MAP}} = \argmax_\theta \pi(\theta \mid x)
  = \argmax_\theta \big[\log L(\theta \mid x) + \log \pi(\theta)\big].
\]
\end{definition}

\begin{remark}
The MAP coincides with the MLE when the prior is uniform. With a
Gaussian prior $\mathcal{N}(0, \tau^2)$, the MAP corresponds to Ridge
regularization:
\[
  \hat{\theta}_{\text{MAP}} = \argmin_\theta \left[
  -\log L(\theta \mid x) + \frac{\norm{\theta}^2}{2\tau^2}\right].
\]
\end{remark}

\begin{proposition}[MAP and Lasso regularization]
If the prior is a Laplace distribution
$\pi(\theta) \propto e^{-\lambda \abs{\theta}}$, the MAP corresponds
to Lasso ($\ell_1$) regularization.
\end{proposition}

% ------------------------------------------------------------
\section{Credible intervals}
% ------------------------------------------------------------

\begin{definition}[Credible interval]
A $100(1-\alpha)\%$ \textbf{credible interval} is an interval $[a, b]$
such that:
\[
  \PP(\theta \in [a, b] \mid x) = 1 - \alpha.
\]
\end{definition}

\begin{definition}[HDI --- Highest Density Interval]
The \textbf{highest density interval} (HDI) is the shortest credible
interval:
\[
  \text{HDI}_{1-\alpha} = \{
  \theta : \pi(\theta \mid x) \geq c_\alpha\},
\]
where $c_\alpha$ is chosen so that
$\PP(\pi(\theta \mid x) \geq c_\alpha \mid x) = 1 - \alpha$.
\end{definition}

\begin{warning}[title=\textbf{Credible vs.\ confidence}]
A $95\%$ credible interval means that $\theta$ belongs to the interval
with probability $0.95$ \textit{conditional on the observed data}. A
$95\%$ confidence interval means that the procedure contains $\theta$
in $95\%$ of \textit{hypothetical} repetitions.
\end{warning}

\begin{example}
For the Bernoulli--Beta model with $n = 100$, $s = 30$,
$\theta \sim \text{Be}(1,1)$, the posterior is $\text{Be}(31, 71)$.

The equal-tailed $95\%$ credible interval is:
\[
  [F^{-1}_{31,71}(0.025),\; F^{-1}_{31,71}(0.975)]
  \approx [0.216,\; 0.397].
\]
\end{example}

% ------------------------------------------------------------
\section{Bayesian vs.\ frequentist comparison}
% ------------------------------------------------------------

\begin{definition}[Frequentist risk]
The frequentist risk of an estimator $\delta(x)$ is:
\[
  R(\theta, \delta) = \E_{x \mid \theta}[L(\theta, \delta(x))].
\]
\end{definition}

\begin{definition}[Admissibility]
An estimator $\delta$ is \textbf{admissible} if there is no estimator
$\delta'$ such that $R(\theta, \delta') \leq R(\theta, \delta)$ for
all $\theta$, with strict inequality for at least one $\theta$.
\end{definition}

\begin{theorem}[Admissibility of Bayes estimators]
Under regularity conditions, every Bayes estimator with finite Bayes
risk is admissible.
\end{theorem}

\begin{theorem}[Complete class]
Under general conditions, the class of Bayes estimators (proper and
limits) forms a \textbf{complete class}: every admissible estimator is
a Bayes estimator (possibly generalized).
\end{theorem}

\begin{example}
\textbf{James--Stein estimator.} For $X \sim \mathcal{N}_p(\theta, I_p)$
with $p \geq 3$, the estimator
$\hat{\theta}_{\text{JS}} = (1 - \frac{p-2}{\norm{X}^2})X$ dominates
the MLE $\hat{\theta} = X$ under quadratic loss. It is a generalized
Bayes estimator.
\end{example}

% ------------------------------------------------------------
\section{Point estimation: examples}
% ------------------------------------------------------------

\begin{example}
\textbf{Poisson--Gamma model.}
With $X_i \sim \text{Poi}(\lambda)$ and $\lambda \sim \text{Ga}(\alpha, \beta)$,
the posterior is $\text{Ga}(\alpha + \sum x_i, \beta + n)$. The point
estimators are:
\begin{align*}
  \hat{\lambda}_{\text{mean}} &= \frac{\alpha + \sum x_i}{\beta + n}, \\
  \hat{\lambda}_{\text{MAP}} &= \frac{\alpha + \sum x_i - 1}{\beta + n}
  \quad (\text{if } \alpha + \sum x_i > 1), \\
  \hat{\lambda}_{\text{median}} &\approx \hat{\lambda}_{\text{mean}}
  - \frac{1}{3(\beta + n)} \quad (\text{approx.\ for large } \alpha').
\end{align*}
\end{example}

\begin{example}
\textbf{Normal--Normal model.}
With $X_i \sim \mathcal{N}(\mu, \sigma^2)$, $\sigma^2$ known, and
$\mu \sim \mathcal{N}(\mu_0, \tau_0^2)$: the posterior is Gaussian
hence symmetric, and all three estimators coincide:
\[
  \hat{\mu}_{\text{mean}} = \hat{\mu}_{\text{MAP}} = \hat{\mu}_{\text{median}}
  = \mu_n = \tau_n^2\left(\frac{\mu_0}{\tau_0^2} + \frac{n\bar{x}}{\sigma^2}\right).
\]
\end{example}

% ------------------------------------------------------------
\section{LINEX loss and asymmetric estimation}
% ------------------------------------------------------------

\begin{definition}[LINEX loss]
The LINEX (\textit{LINear-EXponential}) loss is:
\[
  L(\theta, a) = e^{c(\theta - a)} - c(\theta - a) - 1, \qquad c \neq 0.
\]
For $c > 0$, overestimation is penalized exponentially. For $c < 0$,
underestimation is penalized.
\end{definition}

\begin{theorem}[Bayes estimator under LINEX loss]
The Bayes estimator under LINEX loss is:
\[
  \hat{\theta}_{\text{LINEX}} = -\frac{1}{c}\log \E[e^{-c\theta} \mid x].
\]
\end{theorem}

\begin{proof}
The posterior risk is:
\[
  \rho(a) = \E[e^{c(\theta-a)} \mid x] - c\,\E[\theta - a \mid x] - 1.
\]
Differentiating with respect to $a$ and setting to zero:
\[
  -c\,e^{-ca}\,\E[e^{c\theta} \mid x] + c = 0
  \implies e^{ca} = \E[e^{c\theta} \mid x],
\]
whence $a = \frac{1}{c}\log \E[e^{c\theta} \mid x]$.
\end{proof}

% ------------------------------------------------------------
\section{Bayesian shrinkage}
% ------------------------------------------------------------

\begin{proposition}[Shrinkage toward the prior mean]
In the Normal--Normal model, the posterior mean can be written as:
\[
  \E[\mu \mid x] = (1 - B)\bar{x} + B\mu_0,
  \qquad B = \frac{\sigma^2/n}{\tau_0^2 + \sigma^2/n},
\]
where $B \in (0,1)$ is the \textbf{shrinkage factor}. The Bayesian
estimator is shrunk toward the prior mean $\mu_0$.
\end{proposition}

% ------------------------------------------------------------
\section{Python implementation}
% ------------------------------------------------------------

\begin{codeblock}[title=\textbf{Bayesian estimation and credible intervals}]
\begin{minted}{python}
import numpy as np
from scipy import stats
import matplotlib.pyplot as plt

# Bernoulli data
np.random.seed(42)
n, s = 100, 30

# Posterior Beta(31, 71)
a_post, b_post = 1 + s, 1 + n - s
posterior = stats.beta(a_post, b_post)

# Point estimators
mean_post = posterior.mean()
median_post = posterior.median()
mode_post = (a_post - 1) / (a_post + b_post - 2)

print(f"Posterior mean: {mean_post:.4f}")
print(f"Posterior median: {median_post:.4f}")
print(f"MAP: {mode_post:.4f}")
print(f"MLE: {s/n:.4f}")

# Credible intervals
ci_equi = posterior.ppf([0.025, 0.975])
print(f"Equal-tailed 95% CI: [{ci_equi[0]:.4f}, {ci_equi[1]:.4f}]")

# HDI (via ArviZ)
import arviz as az
samples = posterior.rvs(100000)
hdi = az.hdi(samples, hdi_prob=0.95)
print(f"95% HDI: [{hdi[0]:.4f}, {hdi[1]:.4f}]")
\end{minted}
\end{codeblock}

\begin{codeblock}[title=\textbf{Comparing loss functions}]
\begin{minted}{python}
delta = np.linspace(-3, 3, 500)

fig, axes = plt.subplots(1, 3, figsize=(14, 4))

# Quadratic loss
axes[0].plot(delta, delta**2)
axes[0].set_title("Quadratic loss")
axes[0].set_xlabel(r"$\theta - a$")

# Absolute value loss
axes[1].plot(delta, np.abs(delta))
axes[1].set_title("Absolute value loss")
axes[1].set_xlabel(r"$\theta - a$")

# LINEX loss (c=1 and c=-1)
for c in [1, -1]:
    loss = np.exp(c*delta) - c*delta - 1
    axes[2].plot(delta, loss, label=f"c={c}")
axes[2].set_title("LINEX loss")
axes[2].set_xlabel(r"$\theta - a$")
axes[2].legend()

for ax in axes:
    ax.set_ylabel("Loss")

plt.tight_layout()
plt.savefig("loss_functions.pdf")
\end{minted}
\end{codeblock}

\begin{codeblock}[title=\textbf{Bayesian shrinkage effect}]
\begin{minted}{python}
import pymc as pm
import arviz as az

# Simulation: J groups, simultaneous estimation
np.random.seed(0)
J = 8
true_mu = np.array([-2, -1, 0, 0.5, 1, 1.5, 2, 3])
sigma = 1.0
n_per_group = 5

data = [np.random.normal(mu, sigma, n_per_group)
        for mu in true_mu]
x_bar = np.array([d.mean() for d in data])

# MLE vs Bayes estimators (prior N(0, 4))
tau0 = 2.0
B = (sigma**2 / n_per_group) / (tau0**2 + sigma**2 / n_per_group)
bayes_est = (1 - B) * x_bar + B * 0  # mu_0 = 0

fig, ax = plt.subplots(figsize=(8, 6))
ax.scatter(range(J), true_mu, marker="*", s=150,
           label=r"True $\mu_j$", zorder=3)
ax.scatter(range(J), x_bar, marker="o",
           label=r"MLE $\bar{x}_j$")
ax.scatter(range(J), bayes_est, marker="s",
           label="Bayes (shrinkage)")
ax.axhline(0, color="gray", linestyle=":", alpha=0.5)
ax.set_xlabel("Group $j$")
ax.set_ylabel(r"$\mu_j$")
ax.legend()
ax.set_title("Bayesian shrinkage")
plt.tight_layout()
plt.savefig("shrinkage.pdf")
\end{minted}
\end{codeblock}

% ------------------------------------------------------------
\section{Exercises}
% ------------------------------------------------------------

\begin{exercise}
Prove that the Bayes estimator under the weighted loss
$L(\theta, a) = w(\theta)(\theta - a)^2$ is:
\[
  \hat{\theta}_B = \frac{\E[w(\theta)\theta \mid x]}{\E[w(\theta) \mid x]}.
\]
\end{exercise}

\begin{exercise}
For the Poisson--Gamma model, compute the Bayes estimator under LINEX
loss with $c = 1$ and compare with the posterior mean.
\end{exercise}

\begin{exercise}
Show that in the univariate Gaussian case, the HDI coincides with the
equal-tailed interval.
\end{exercise}

\begin{exercise}
Let $X \sim \mathcal{N}_p(\theta, I_p)$ with
$\theta \sim \mathcal{N}_p(0, \tau^2 I_p)$.
\begin{enumerate}
  \item Compute the Bayes estimator under quadratic loss.
  \item Show that it is a shrinkage estimator of the form $c \cdot X$
        with $c \in (0,1)$.
  \item Compare with the James--Stein estimator.
\end{enumerate}
\end{exercise}

\begin{exercise}
\textbf{(Numerical)} Simulate $n = 50$ observations from
$\mathcal{N}(3, 4)$. Compare the Bayesian estimators of $\mu$ (mean,
median, MAP) with the MLE for different priors:
$\mathcal{N}(0, 100)$ (vague), $\mathcal{N}(0, 1)$ (informative,
centered), $\mathcal{N}(10, 1)$ (informative, off-center).
\end{exercise}
